\begin{definition}
Let $n\ge2$ and $p>1$.  The predicate
$\mathrm{CitedHolderOscillationSolution}(n,p)$ is part~(b) of the cited result
(C4) of the Note for the formalization (De
Giorgi--Ladyzhenskaya--Ural'tseva H\"older estimates;
[G, chapter VII, Theorems 2.1--2.2], [GG2], [LU, chapter 4]):

\emph{for all $\mu,\gamma_0,\dots,\gamma_4>0$
and $M\ge0$ there are $\alpha\in(0,1]$ and $d>0$ such that, for every
$\varepsilon\in[0,1]$: if $\Omega$ is open, $(A,B)$ satisfies
(Q1) of the thesis--(Q5) of the thesis on $\Omega$ with these parameters
(\noderef{EllipticStructureQ12}, \noderef{EllipticStructureQ345}), $u$ is a weak
solution of (1.1) of the thesis in $\Omega$ (\noderef{IsWeakSolutionOn}) with
$|u|\le M$ a.e.\ on $\Omega$, and $\overline B_R(x_0)\subseteq\Omega$, then for
$0<\rho\le R$}
\[
\operatorname*{osc}_{B_\rho}u\;\le\;4^\alpha\max\Bigl\{\operatorname*{osc}_{B_R}u,\;dR\Bigr\}\Bigl(\frac\rho R\Bigr)^\alpha .
\]

This is (5.1) of the thesis.  The constants $\alpha$ and $d$
are bound after $n,p$, the structure constants and $M=\|u\|_{L^\infty}$, and
before $\varepsilon$, $\Omega$, the operator pair and $u$.  That is exactly the
dependence the Note allows: at l.~263 it says the constants ``may be allowed to
depend also on $M=\|u\|_{L^\infty}$'', naming $M$ and nothing else, so
$\varepsilon$ is not among the quantities $\alpha$ and $d$ may depend on and
one pair of constants must serve the whole range $0\le\varepsilon\le1$.

\medskip
\noindent\textbf{Why this is a predicate of its own and not a field of a
two-part bundle, and what makes it true.}  The Note at l.~263 labels the two
halves of (C4) as ``(a) A bounded local minimum \ldots'' and ``(b) a bounded
weak solution \ldots'', and they share no parameter beyond $n$ and $p$.  Keeping
them as two fields of one structure would mean that a consumer of (a) assumes
(b) as well, so a failure of either would empty the targets of both; the part
(a) predicate is therefore \noderef{CitedHolderOscillationMinimum} and this one
stands alone.

The hypothesis that makes this statement true, and whose absence would make it
false, is the Carath\'eodory pair now carried by
\noderef{EllipticStructureQ12} and \noderef{EllipticStructureQ345}, together
with the integrability clause of \noderef{IsWeakSolutionOn}.  Without them a
coefficient field may be taken non-measurable in $x$ --- for instance
$A(x,z,h)=(1+\mathbf1_V(x))(\varepsilon+|h|^2)^{(p-2)/2}h$ with $V$ a
non-measurable set --- in which case the integrand of (1.2) of the thesis is
non-measurable, the weak-formulation identity holds for \emph{every}
$u\in W^{1,p}_{\mathrm{loc}}(\Omega)$ by the junk value of the integral, and no
uniform oscillation decay can hold: the functions $u_k(x)=\sin(kx_1)$ would all
be ``weak solutions'' while $\operatorname*{osc}_{B_\rho}u_k$ does not decay
uniformly in $k$.  The sources the Note cites at l.~263 assume Carath\'eodory
data, so the clauses restore the literature's hypotheses, which l.~258 requires,
rather than adding to them.

\medskip
\noindent\textbf{Why $\varepsilon$ is bound after the two constants, and not
before.}  This is not a strengthening beyond the literature; it is what the
literature states and what \S4 needs.  The De Giorgi estimates are proved from
the growth condition (F1) of the thesis (resp.\ (Q1) of the thesis--(Q5) of the thesis),
whose weight is sandwiched uniformly over the range in question:
\[
|h|^{p}\;\le\;\bigl(\varepsilon+|h|^{2}\bigr)^{\frac p2}\;\le\;\bigl(1+|h|^{2}\bigr)^{\frac p2}
\qquad(0\le\varepsilon\le1),
\]
so the constants the iteration produces depend on the two ends of that
sandwich, not on where $\varepsilon$ sits inside it.  The ordering is also
forced by the use: Lemma~Lemma~4.1 of the thesis's own list of what $\sigma_2$ and $R_1$
depend on (thesis l.~1163--1164) is $n,p,\lambda,\Lambda,a,\nu,M$ and
$d(\Omega',\partial\Omega)$, with no $\varepsilon$, while $\sigma_0$ and
$\sigma_1$ carry $\sigma(R+C_1R^{\alpha})$ and $C_1R^{\alpha}$ and so make
$\sigma_2$ inherit $R^{\alpha}$.  With $\varepsilon$ bound before the
existential, $\alpha$ would be a function $\alpha(\varepsilon)$ and a single
$\sigma_2$ would have to dominate $R^{\alpha(\varepsilon)}$ for every
$\varepsilon$; since $R^{\alpha}\to1$ as $\alpha\downarrow0$ for each fixed
$R\in(0,1)$, the pointwise supremum of that family is identically $1$ on
$(0,1)$, which no function with $\sigma_2(0)=0$ and $\sigma_2$ increasing can
dominate.  So the earlier ordering made Lemma~Lemma~4.1 of the thesis unprovable rather
than merely harder.  All
oscillations are essential oscillations (\noderef{EssOscOn}); the hypothesis
$|u|\le M$ a.e.\ makes them well defined.  This result is used as an explicit
hypothesis and is never introduced as an axiom.
\end{definition}
